\scanpage{133}
% archive sheet number: 132

qui prolongent $\varphi_H,\varphi_{H'}$.

$\deleted{\unsure{T}}L_P$: Ens des formes \uline{lin.} sur $P$ qui prolongent $T_B$

On a des flèches canoniques

\diagram{\begin{tikzcd}[ampersand replacement=\&,row sep=tiny,column sep=small]
\deleted{\unsure{Q_P|\underline{\Phi}_P}}\ \deleted{\underline{\Phi}_P}\arrow[r,"{\deleted{\to}}" description] \& \deleted{\underline{\Phi}_B}\arrow[r,"{\deleted{\to}}" description] \& \deleted{\unsure{\underline{\Phi}|\underline{\Phi}_B}}
\end{tikzcd}}

\diagram{\begin{tikzcd}[ampersand replacement=\&,row sep=small,column sep=small]
\underline{Q}_P\arrow[r,"\alpha"]\arrow[dr,"\delta",swap]\arrow[dr,phantom,shift right=2ex,"{\deleted{\unsure{\varphi}}}"'] \& \underline{Q}_B\arrow[r,"\beta"] \& {\deleted{\unsure{\underline{\Phi}_B}}\,\underline{\Psi}_B}\arrow[r,"\gamma"] \& \underline{\Phi}_B\arrow[dll,"\zeta"]\\
\& \underline L_P \& \deleted{\unsure{\varphi_P\to\underline{\Phi}_B}} \&
\end{tikzcd}}

que voici : $\alpha$ est la restriction $\alpha(\deleted{\unsure{q}}\Delta_P)=\deleted{\unsure{q}}\Delta\circ\alpha$,
$\beta$ associe à tte forme $\Delta_B\in Q_B$ la
forme bil. sym. associée \inserted{$\psi_B$}, enfin
on définit $\gamma$ par la formule

\[
\underset{\varphi_B}{\underbrace{\gamma(\psi)}_{\parallel\,\mathrm{déf}}}
\deleted{\unsure{=}}+\psi_B=\deleted{\unsure{\{}}
\bigl((u,v)\longmapsto T_B(u)T_B(v)\bigr)
\]

Enfin, on définit $\delta$ par

\[
\delta(\Delta_P)=\{u\longmapsto\varphi_{\Delta_P}(u,1_P)\}
\]

% slanted
\margin{On définit $\zeta:\underline{\Phi}_B\to\underline L_P$ ainsi : $\underline L_P$ s'identifie \inserted{à l'}\deleted{\unsure{au|avec|aux}}ens des formes \uline{bilin.} sur $H\times H'$ telles que leurs restrictions à $1\times H'$, et à $H\times1$ sont données par $T_{H'}$ et $T_H$. à $\varphi_B$ on associe la forme bil. $\zeta(\varphi_B)\in\unsure{*|:}\,(u,v)\longmapsto\varphi_B(u,v)$. C'est encore un hom de torseurs sous $(L\otimes L')^\vee$, donc iso.}

Ceci dit, les quatre applications
$\alpha,\beta,\gamma,\delta$ sont \uline{bijectives}.

\uline{Remarques} 1) Pour définir les cinq
ensembles $\underline Q_P,\underline L_P,\underline Q_B,\underline\Psi_B,\underline\Psi_B$, \deleted{et}
les applications $\alpha,\beta,\gamma,\delta$, on ne
utilise sur $H,H'$ que la donnée
de $\Delta_H,\Delta_{H'}$ et $1_H,1_{H'}$ ($\unsure{\longrightarrow|\Longrightarrow}T_H,T_{H'}$ \unsure{etc|cf.})
satisfaisant $\Delta_H(1_H)=\Delta_{H'}(1_{H'})=1$.

